\scanpage{75}\mspage{1}
Soit $E$ une courbe elliptique -- on travaille sur un corps alg. clos $k$ pour simplifier le propos, mais la seule hypothèse essentielle est que la base soit à car. résiduelles $\ne2$. On désigne par \deleted{\unsure{$Z$|$J$|$\Sigma$}}
\[
J={}_2E(k)^\ast \tag*{(1)}
\]
l'ens des pts de $E$ d'ordre 2, de sorte que l'on a
\[
E/(\pm\mathrm{id}_E)\simeq\Sigma_J\qquad
\text{\deleted{\unsure{quadrique de|quotient de|spectre de}} \inserted{droite projective} $J$-marquée.} \tag*{(2)}
\]
\[
\Sigma_J\simeq\mathbb P(V_J(k))\qquad
0\mathbin{-} V_J(k)\to k^J\xrightarrow{\mathrm{somme}} k\to 0 \tag*{(3)}
\]

La donnée de $E$ est équivalente à celle de $J$ i.e. de $\Sigma_J$ muni du diviseur étale $D_J$ de degré 3, plus \unsure{le|les} point
\[
t\in\Sigma_J^\ast=\Sigma_J\setminus D_J \tag*{(4)}
\]
image de l'élément $0_E$, plus une droite $T$ \inserted{sur $k$} (savoir $T_{E,0}$), munie d'un isomorphisme
\[
T^{\otimes2}\xrightarrow{\sim}T_{\Sigma,t} \tag*{(5)}
\]
(\deleted{\unsure{Réinterprétation|Réinterpréter|isomorphisme|isomorphismes}} \inserted{équivalence} de catégories-groupoïdes -- valable en fait sur base $S$ quelconque à car. rés. $\ne2$, et compatible avec changement de base \ldots).

1 \quad Notons qu'on a un isomorphisme canonique
\[
T_\Sigma\simeq\underline O_\Sigma(2)(\underline{\unsure{\omega|w}}) \tag*{(6)}
\]
où $\underline\omega=\underline\omega(J)$ est l'ens. des deux ordres circulaires sur $J$, et où $L(\underline\omega)$ (pour un objet tel que faisceau abélien etc) désigne $L\wedge_{\{\pm1\}}\underline\omega$.
% archive sheet number: 74
